\label{sec:interp}

A natural follow-up to the empirical claims in Sections~\ref{sec:results-curve}--\ref{sec:diff-draft-refill}:
\emph{do the AR head and the diffusion head use the same internal
features inside the shared backbone, or do they specialise?} Phase~P
answers this with a mechanistic-interpretability layer built on top
of the F10 ($305$\,M, $d{=}1024$, $L{=}20$) checkpoint --- no retrain,
no architecture change, $\sim\$3$ of compute. We adapt the
SAE-for-diffusion methodology of DLM-Scope~\citep{wang2026dlmscope}
(which studies \emph{pure} diffusion LMs) to the composite setting,
where the new question is whether the two heads sharing one backbone
specialise or share features.

\subsection{Setup}

We adopt nnsight \citep{nnsight2024} to wrap the custom
\code{CompositeLM} (Sfumato is not a HuggingFace model). We train
TopK sparse autoencoders \citep{topk_sae_2024} on F10's residual
stream:
\begin{itemize}
  \setlength\itemsep{2pt}
  \item Ten per-layer SAEs on the backbone, layers $6$--$15$
        (the divergence zone identified in Phase~P.0 below),
        AR-mode activations.
  \item Two pre-head SAEs on \code{ln\_f} output, one trained on
        AR-mode activations and one on diff-mode activations. These
        let us directly compare what each head ``sees''.
\end{itemize}
All SAEs are $k{=}64$ TopK with $16{,}384$ features
(expansion $16\times$ over $d{=}1024$), trained for $8{,}000$ steps
on $200$\,M FineWeb-Edu tokens on a single A40 ($\sim$3\,h, $\sim\$2$).
All checkpoints meet the standard reconstruction target
(delta-loss $< 0.10$ on held-out tokens).

\subsection{P.0 — raw activation divergence by layer}

Before training any SAEs we compute, for a single GSM8K prompt run
through F10 in both modes, the cosine similarity of the residual
stream at each layer:

\begin{center}
\begin{tabular}{lrr}
\toprule
layer & $\cos(\mathrm{AR}, \mathrm{diff})$ & reading \\
\midrule
0 & $0.987$ & tok+pos emb, modes identical \\
5 & $0.880$ & early shared processing \\
10 & $0.650$ & divergence zone \\
14 & $0.510$ & maximum specialisation \\
19 & $0.543$ & slight rebound at \code{ln\_f} \\
\bottomrule
\end{tabular}
\end{center}

The drop is concentrated in layers $6$--$13$. We focus SAE training
there.

\subsection{P.2 — cross-head SAE feature overlap (headline)}

We train two TopK SAEs at the same residual position (post-\code{ln\_f}),
one on AR-mode activations and one on diff-mode activations, then
compute the cosine between every AR decoder column and every diff
decoder column:

\begin{center}
\begin{tabular}{lr}
\toprule
metric & value \\
\midrule
mean best-cosine (AR$\to$diff) & $0.169$ \\
median & $0.129$ \\
maximum & $0.904$ \\
\% AR features with diff sibling $\geq 0.5$ & $1.84$\,\% \\
\% AR features with diff sibling $\geq 0.7$ & $0.42$\,\% \\
\% AR features with diff sibling $\geq 0.9$ & $0.01$\,\% (2 of $16{,}384$) \\
\bottomrule
\end{tabular}
\end{center}

\paragraph{P.2c --- the control that anchors the number.}
A mean best-cosine of $0.169$ is only meaningful against two
baselines, because TopK SAEs have no canonical basis: two SAEs need
not align even on identical data. We measure both.
\emph{(i) Chance floor.} The mean best-cosine of the AR decoder to a
random unit dictionary of the same size ($16{,}384$ columns in
$\mathbb{R}^{1024}$) is $0.124$ --- in this dimensionality, best-of-many
cosine is non-trivial by construction.
\emph{(ii) Same-mode null.} We retrain a \emph{second} AR-mode SAE at
a different seed and compute AR-seed0 $\to$ AR-seed1: $0.286$.

\begin{center}
\begin{tabular}{lr}
\toprule
comparison & mean best-cosine \\
\midrule
AR-seed0 $\to$ AR-seed1 (same mode, diff.\ seed) & $0.286$ \\
\textbf{AR $\to$ diff (cross-mode)} & $\mathbf{0.169}$ \\
AR $\to$ random dictionary (chance floor) & $0.124$ \\
\bottomrule
\end{tabular}
\end{center}

\noindent The cross-mode alignment ($0.169$) sits \emph{below} the
same-mode baseline ($0.286$) and only $+0.045$ above chance, whereas
the same-mode baseline is $+0.162$ above chance. Going from same-mode
to cross-mode therefore loses $\sim\!72\%$ of the above-chance feature
alignment ($1 - 0.045/0.162$). This is the rigorous statement of the
result: \textbf{the two heads share substantially less feature
structure than two independently-trained SAEs of the \emph{same} head
do.} (The seed-1 SAE was trained for $2{,}000$ steps vs.\ $8{,}000$
for the originals, so $0.286$ is a conservative --- low --- estimate
of same-mode alignment; a matched run would only widen the gap.) The
pre-committed ``$<\!0.4$ = specialise'' threshold is thus empirically
anchored: the relevant comparison is the $0.286$ same-mode baseline,
not an arbitrary cut.

We therefore avoid the over-strong word ``disjoint'': cross-mode
cosine is above chance, so a small shared ``bridge'' subspace exists
(consistent with the $\sim\!20$ high-cosine matched pairs in P.3).
The defensible claim is \textbf{substantial specialisation}, not
orthogonality.

\textbf{The composite at $305$\,M does not converge to ``one
representation, two readouts''.} The two heads operate on largely
distinct sparse feature bases inside the same backbone weights.

\subsection{P.2b — specialisation by depth}
\label{sec:p2b}

The P.2 measurement above is at \code{ln\_f} only (the pre-head
residual). To test whether specialisation is a pre-head artefact
(arising from the explicit \code{head\_diff\_proj} projection) or a
backbone-wide property, we train symmetric AR-mode and diff-mode
SAEs at every layer in the divergence zone (blocks $6$--$15$, ten
blocks each side) and compute cross-head best-cosine at each depth.

\begin{center}
\begin{tabular}{lrrrr}
\toprule
layer & mean cos & median & max & \% sibling $\geq 0.7$ \\
\midrule
block\_6  & $0.162$ & $0.128$ & $0.924$ & $0.2$\,\% \\
block\_7  & $0.159$ & $0.128$ & $0.905$ & $0.1$\,\% \\
block\_8  & $0.161$ & $0.128$ & $0.914$ & $0.1$\,\% \\
block\_9  & $0.162$ & $0.127$ & $0.924$ & $0.1$\,\% \\
block\_10 & $0.161$ & $0.127$ & $0.932$ & $0.1$\,\% \\
block\_11 & $0.159$ & $0.126$ & $0.913$ & $0.1$\,\% \\
block\_12 & $0.159$ & $0.126$ & $0.951$ & $0.1$\,\% \\
block\_13 & $0.158$ & $0.126$ & $0.845$ & $0.1$\,\% \\
block\_14 & $0.158$ & $0.127$ & $0.845$ & $0.1$\,\% \\
block\_15 & $0.157$ & $0.127$ & $0.904$ & $0.2$\,\% \\
\midrule
\code{ln\_f} & $0.169$ & $0.129$ & $0.904$ & $0.4$\,\% \\
\bottomrule
\end{tabular}
\end{center}

\textbf{Specialisation is uniform-by-depth.} Block-level mean cosine
sits in the tight band $[0.157, 0.162]$ across blocks $6$--$15$ (range
$0.005$, coefficient of variation $\sim 1$\,\%). The ``modes
specialise'' finding is a backbone-wide property, not a pre-head
phenomenon. Notably, \code{ln\_f} ($0.169$) is slightly \emph{less}
specialised than internal blocks --- opposite to what we'd expect if
the explicit \code{head\_diff\_proj} layer caused the specialisation.
The bridge-feature tail (cosine $\geq 0.7$) accounts for under
$0.5$\,\% of features at every depth. The original P.2 measurement
at \code{ln\_f} was, if anything, a slight under-statement of how
specialised the circuits are internally.

\subsection{P.3 — per-prompt confirmation across 12 prompts}

We rerun on 4 GSM8K, 4 prose, and 4 explicit math-mode prompts;
for each prompt we extract the top-$8$ firing AR features and
top-$8$ firing diff features and check the index overlap.

\textbf{Across all 12 prompts the overlap is exactly $0$.} The two
modes fire \emph{completely different} feature indices when
processing the same prompt. Mean token-level cosine ranges from
$0.439$ (M1 simple arithmetic) to $0.817$ (M2 chained $x$ definition).

The $20$ ``bridge features'' (P.2 top-$20$ matched pairs with cosine
$0.85$--$0.90$) appear in $1$--$2$ top-$8$ slots per prompt and
fire on broadly content-anchor tokens. These are candidate
mode-switch routing features for the Phase K recipe.

\subsection{P.4 — causal validation via steering}

We pick a top AR-mode-specific feature (\#9304, P.3-identified) and
add $+k$\,$\times$ its SAE decoder column to F10's \code{ln\_f}
output during greedy generation. Compared to unsteered baseline on
the same prompt:

\begin{itemize}
  \setlength\itemsep{2pt}
  \item Prompt where the feature naturally fires (Janet ducks):
        $100$\,\% token-equal generation at $k\!=\!4$ and $k\!=\!8$.
  \item Prompt where it does not (Roman Empire history): generation
        diverges by $43\!-\!64$\,\%, breaking the AR-mode repetition
        loop the baseline exhibits.
  \item Math-chain prompt: $48$\,\%--$77$\,\% divergence depending on
        mode and $k$.
\end{itemize}

The SAE-extracted feature directions are not just correlational ---
they have measurable causal influence on F10 generation. The
content-conditional behaviour (large effect only on prompts where
the feature already participates) is the expected pattern for a
learned direction.

\subsection{What this opens}

\paragraph{Distillation candidate.} If AR and diff feature bases are
substantially specialised, F10's backbone is effectively two specialised
subnetworks sharing only the bottom-most layers and a small
``bridge'' subspace. A distillation into AR-only and diff-only
models with shared embeddings could recover most of F10's capability
at $\sim$half the parameter count. Worth a paper-D scaling experiment.

\paragraph{Phase K bridge features.} The $\sim$20 high-cosine matched
pairs from P.2 are the natural candidates for the ``mode-switch
routing'' features that K.2's per-token cross-mode recipe
(\textsection\ref{sec:diff-draft-refill}) implicitly uses. They fire
on universal content tokens across all prompt types tested. Future
work: confirm they are the actual mechanism via attribution graphs.

\paragraph{Honest negative.} The interpretability story does not
rescue F10's free-run accuracy (still $\sim$6\,\% on GSM8K-dev).
Mechanistic insight here is about \emph{what the model is computing},
not about closing the capability gap. We document this so readers
do not over-claim the interpretability layer.

\paragraph{Total interpretability cost.} $\sim\$3$ on one A40 for
$\sim$3\,h ($12$ SAEs $+$ overlap analysis). Reproducible via the
recipe at \url{https://github.com/eren23/sfumato/tree/main/e5/interp}
and the SAE weights at
\url{https://huggingface.co/eren23/sfumato-composite-ckpts/tree/main/interp/saes}.
